\documentclass[10pt,letterpaper]{article}
\usepackage{cornell-notes}
\usepackage{mathtools}

\title{Chapter 15: Modeling of Data}
\author{}
\date{}

\setDocTitle{Chapter 15: Modeling of Data}
\setDocSubtitle{Numerical Methods}
\setDocVersion{v1.0}
\setDocStatus{Study Companion}
\setDocOwner{Numerical Methods Cornell Notes}
\setDocAuthor{}
\setDocDate{\today}
\setDocSummary{Cornell-format study and review notes for Chapter 15: Modeling of Data.}

\setCornellCollection{Numerical Methods Cornell Notes}
\setCornellUnitType{Chapter}
\setCornellUnitNumber{15}
\setCornellUnitTitle{Modeling of Data}
\setCornellCompanionLabel{Study and review companion}

\hypersetup{pdftitle={Chapter 15: Modeling of Data},pdfsubject={Cornell notes},pdfauthor={}}

\begin{document}
\maketitle

\section*{Chapter Overview}
This chapter organizes the mathematical and computational ideas behind modeling of data. The notes preserve the numbered section sequence while emphasizing method selection, explicit assumptions, numerical reliability, cost, and verification.

\section*{Learning Objectives}
Explain the governing problem class; compare Least Squares as a Maximum Likelihood Estimator, Fitting Data to a Straight Line, Straight-Line Data with Errors in Both Coordinates, and General Linear Least Squares; design defensible stopping and verification checks; distinguish problem conditioning from algorithmic stability.

\section*{Prerequisites}
Probability, linear algebra, optimization, likelihood, and matrix factorization.

\section*{Operational Workflow}
Specify observation and error models; build the design or covariance structure; fit with a stable solver; inspect residuals and identifiability; quantify uncertainty; test predictions out of sample.

\subsection*{Formula and notation anchor}
\[
\widehat{\beta}=\arg\min_{\beta}(y-X\beta)^T W(y-X\beta)
\]
Normal equations square the condition number; QR or SVD is generally preferred for a stable least-squares implementation.

\begin{warnbox}{Numerical warning}
Rank deficiency, overfitting, correlated errors, unjustified Gaussian assumptions, parameter nonidentifiability, poor MCMC mixing, kernel hyperparameter confounding, and data leakage.
\end{warnbox}

\section{15.0 Introduction}
\begin{cornell}
\noterow{What is the central idea?}{Data modeling links a stochastic error model to parameter estimation, diagnostics, uncertainty, prediction, and model criticism.}
\noterow{How should it be used?}{For Introduction, begin from explicit inputs, outputs, preconditions, and representation choices.  Specify observation and error models; build the design or covariance structure; fit with a stable solver; inspect residuals and identifiability; quantify uncertainty; test predictions out of sample.}
\noterow{Cost and reliability}{Analyze arithmetic work, storage, data movement, and convergence separately.  Relevant failure pressures include rank deficiency, overfitting, correlated errors, unjustified gaussian assumptions, parameter nonidentifiability, poor mcmc mixing, kernel hyperparameter confounding, and data leakage.  Use scaled tolerances and report both success criteria and diagnostic evidence.}
\noterow{Implementation checklist}{\begin{itemize}
\item Validate dimensions, domains, ordering, and structural assumptions.
\item Guard small divisors, invalid states, overflow, underflow, and nonfinite values.
\item Make mutation, workspace, indexing, and termination behavior explicit.
\item Test a simple exact case, a difficult valid case, a boundary case, and an expected failure.
\end{itemize}}
\end{cornell}
\begin{summarybox}{Section summary}
Data modeling links a stochastic error model to parameter estimation, diagnostics, uncertainty, prediction, and model criticism.  Select this technique only when its assumptions match the data, and accept a result only after an independent residual, error estimate, invariant, or refinement check.
\end{summarybox}

\section{15.1 Least Squares as a Maximum Likelihood Estimator}
\begin{cornell}
\noterow{What is the central idea?}{Under independent Gaussian errors with known scales, minimizing weighted squared residuals is equivalent to maximizing likelihood.}
\noterow{How should it be used?}{For Least Squares as a Maximum Likelihood Estimator, begin from explicit inputs, outputs, preconditions, and representation choices.  Specify observation and error models; build the design or covariance structure; fit with a stable solver; inspect residuals and identifiability; quantify uncertainty; test predictions out of sample.}
\noterow{Quantitative anchor}{\[
\min_\beta\|W^{1/2}(y-X\beta)\|_2^2
\]
State the dimensions, scaling, and validity assumptions before using this relation.  Check the resulting residual or invariant rather than trusting an iteration count alone.}
\noterow{Cost and reliability}{Analyze arithmetic work, storage, data movement, and convergence separately.  Relevant failure pressures include rank deficiency, overfitting, correlated errors, unjustified gaussian assumptions, parameter nonidentifiability, poor mcmc mixing, kernel hyperparameter confounding, and data leakage.  Use scaled tolerances and report both success criteria and diagnostic evidence.}
\noterow{Implementation checklist}{\begin{itemize}
\item Validate dimensions, domains, ordering, and structural assumptions.
\item Guard small divisors, invalid states, overflow, underflow, and nonfinite values.
\item Make mutation, workspace, indexing, and termination behavior explicit.
\item Test a simple exact case, a difficult valid case, a boundary case, and an expected failure.
\end{itemize}}
\end{cornell}
\begin{summarybox}{Section summary}
Under independent Gaussian errors with known scales, minimizing weighted squared residuals is equivalent to maximizing likelihood.  Select this technique only when its assumptions match the data, and accept a result only after an independent residual, error estimate, invariant, or refinement check.
\end{summarybox}

\section{15.2 Fitting Data to a Straight Line}
\begin{cornell}
\noterow{What is the central idea?}{Straight-line fitting estimates slope and intercept from centered sums, with uncertainty determined by error assumptions and leverage.}
\noterow{How should it be used?}{For Fitting Data to a Straight Line, begin from explicit inputs, outputs, preconditions, and representation choices.  Specify observation and error models; build the design or covariance structure; fit with a stable solver; inspect residuals and identifiability; quantify uncertainty; test predictions out of sample.}
\noterow{Quantitative lens}{Identify the governing equation, recurrence, probability law, or geometric predicate for Fitting Data to a Straight Line.  Define every variable and scale; then derive a residual, error estimate, or invariant that can be checked independently.}
\noterow{Cost and reliability}{Analyze arithmetic work, storage, data movement, and convergence separately.  Relevant failure pressures include rank deficiency, overfitting, correlated errors, unjustified gaussian assumptions, parameter nonidentifiability, poor mcmc mixing, kernel hyperparameter confounding, and data leakage.  Use scaled tolerances and report both success criteria and diagnostic evidence.}
\noterow{Implementation checklist}{\begin{itemize}
\item Validate dimensions, domains, ordering, and structural assumptions.
\item Guard small divisors, invalid states, overflow, underflow, and nonfinite values.
\item Make mutation, workspace, indexing, and termination behavior explicit.
\item Test a simple exact case, a difficult valid case, a boundary case, and an expected failure.
\end{itemize}}
\end{cornell}
\begin{summarybox}{Section summary}
Straight-line fitting estimates slope and intercept from centered sums, with uncertainty determined by error assumptions and leverage.  Select this technique only when its assumptions match the data, and accept a result only after an independent residual, error estimate, invariant, or refinement check.
\end{summarybox}

\section{15.3 Straight-Line Data with Errors in Both Coordinates}
\begin{cornell}
\noterow{What is the central idea?}{When both coordinates are uncertain, ordinary vertical residuals are inadequate and the likelihood must account for two-axis error geometry.}
\noterow{How should it be used?}{For Straight-Line Data with Errors in Both Coordinates, begin from explicit inputs, outputs, preconditions, and representation choices.  Specify observation and error models; build the design or covariance structure; fit with a stable solver; inspect residuals and identifiability; quantify uncertainty; test predictions out of sample.}
\noterow{Quantitative lens}{Identify the governing equation, recurrence, probability law, or geometric predicate for Straight-Line Data with Errors in Both Coordinates.  Define every variable and scale; then derive a residual, error estimate, or invariant that can be checked independently.}
\noterow{Cost and reliability}{Analyze arithmetic work, storage, data movement, and convergence separately.  Relevant failure pressures include rank deficiency, overfitting, correlated errors, unjustified gaussian assumptions, parameter nonidentifiability, poor mcmc mixing, kernel hyperparameter confounding, and data leakage.  Use scaled tolerances and report both success criteria and diagnostic evidence.}
\noterow{Implementation checklist}{\begin{itemize}
\item Validate dimensions, domains, ordering, and structural assumptions.
\item Guard small divisors, invalid states, overflow, underflow, and nonfinite values.
\item Make mutation, workspace, indexing, and termination behavior explicit.
\item Test a simple exact case, a difficult valid case, a boundary case, and an expected failure.
\end{itemize}}
\end{cornell}
\begin{summarybox}{Section summary}
When both coordinates are uncertain, ordinary vertical residuals are inadequate and the likelihood must account for two-axis error geometry.  Select this technique only when its assumptions match the data, and accept a result only after an independent residual, error estimate, invariant, or refinement check.
\end{summarybox}

\section{15.4 General Linear Least Squares}
\begin{cornell}
\noterow{What is the central idea?}{A model linear in its parameters becomes a design-matrix solve, preferably via QR or SVD when rank or conditioning is uncertain.}
\noterow{How should it be used?}{For General Linear Least Squares, begin from explicit inputs, outputs, preconditions, and representation choices.  Specify observation and error models; build the design or covariance structure; fit with a stable solver; inspect residuals and identifiability; quantify uncertainty; test predictions out of sample.}
\noterow{Quantitative anchor}{\[
\min_\beta\|W^{1/2}(y-X\beta)\|_2^2
\]
State the dimensions, scaling, and validity assumptions before using this relation.  Check the resulting residual or invariant rather than trusting an iteration count alone.}
\noterow{Cost and reliability}{Analyze arithmetic work, storage, data movement, and convergence separately.  Relevant failure pressures include rank deficiency, overfitting, correlated errors, unjustified gaussian assumptions, parameter nonidentifiability, poor mcmc mixing, kernel hyperparameter confounding, and data leakage.  Use scaled tolerances and report both success criteria and diagnostic evidence.}
\noterow{Implementation checklist}{\begin{itemize}
\item Validate dimensions, domains, ordering, and structural assumptions.
\item Guard small divisors, invalid states, overflow, underflow, and nonfinite values.
\item Make mutation, workspace, indexing, and termination behavior explicit.
\item Test a simple exact case, a difficult valid case, a boundary case, and an expected failure.
\end{itemize}}
\end{cornell}
\begin{summarybox}{Section summary}
A model linear in its parameters becomes a design-matrix solve, preferably via QR or SVD when rank or conditioning is uncertain.  Select this technique only when its assumptions match the data, and accept a result only after an independent residual, error estimate, invariant, or refinement check.
\end{summarybox}

\section{15.5 Nonlinear Models}
\begin{cornell}
\noterow{What is the central idea?}{Nonlinear fitting iteratively linearizes residuals and uses damping or trust control when a full parameter correction is unsafe.}
\noterow{How should it be used?}{For Nonlinear Models, begin from explicit inputs, outputs, preconditions, and representation choices.  Specify observation and error models; build the design or covariance structure; fit with a stable solver; inspect residuals and identifiability; quantify uncertainty; test predictions out of sample.}
\noterow{Quantitative lens}{Identify the governing equation, recurrence, probability law, or geometric predicate for Nonlinear Models.  Define every variable and scale; then derive a residual, error estimate, or invariant that can be checked independently.}
\noterow{Cost and reliability}{Analyze arithmetic work, storage, data movement, and convergence separately.  Relevant failure pressures include rank deficiency, overfitting, correlated errors, unjustified gaussian assumptions, parameter nonidentifiability, poor mcmc mixing, kernel hyperparameter confounding, and data leakage.  Use scaled tolerances and report both success criteria and diagnostic evidence.}
\noterow{Implementation checklist}{\begin{itemize}
\item Validate dimensions, domains, ordering, and structural assumptions.
\item Guard small divisors, invalid states, overflow, underflow, and nonfinite values.
\item Make mutation, workspace, indexing, and termination behavior explicit.
\item Test a simple exact case, a difficult valid case, a boundary case, and an expected failure.
\end{itemize}}
\end{cornell}
\begin{summarybox}{Section summary}
Nonlinear fitting iteratively linearizes residuals and uses damping or trust control when a full parameter correction is unsafe.  Select this technique only when its assumptions match the data, and accept a result only after an independent residual, error estimate, invariant, or refinement check.
\end{summarybox}

\section{15.6 Confidence Limits on Estimated Model Parameters}
\begin{cornell}
\noterow{What is the central idea?}{Local covariance, profile methods, resampling, or posterior summaries quantify parameter uncertainty under increasingly flexible assumptions.}
\noterow{How should it be used?}{For Confidence Limits on Estimated Model Parameters, begin from explicit inputs, outputs, preconditions, and representation choices.  Specify observation and error models; build the design or covariance structure; fit with a stable solver; inspect residuals and identifiability; quantify uncertainty; test predictions out of sample.}
\noterow{Quantitative lens}{Identify the governing equation, recurrence, probability law, or geometric predicate for Confidence Limits on Estimated Model Parameters.  Define every variable and scale; then derive a residual, error estimate, or invariant that can be checked independently.}
\noterow{Cost and reliability}{Analyze arithmetic work, storage, data movement, and convergence separately.  Relevant failure pressures include rank deficiency, overfitting, correlated errors, unjustified gaussian assumptions, parameter nonidentifiability, poor mcmc mixing, kernel hyperparameter confounding, and data leakage.  Use scaled tolerances and report both success criteria and diagnostic evidence.}
\noterow{Implementation checklist}{\begin{itemize}
\item Validate dimensions, domains, ordering, and structural assumptions.
\item Guard small divisors, invalid states, overflow, underflow, and nonfinite values.
\item Make mutation, workspace, indexing, and termination behavior explicit.
\item Test a simple exact case, a difficult valid case, a boundary case, and an expected failure.
\end{itemize}}
\end{cornell}
\begin{summarybox}{Section summary}
Local covariance, profile methods, resampling, or posterior summaries quantify parameter uncertainty under increasingly flexible assumptions.  Select this technique only when its assumptions match the data, and accept a result only after an independent residual, error estimate, invariant, or refinement check.
\end{summarybox}

\section{15.7 Robust Estimation}
\begin{cornell}
\noterow{What is the central idea?}{Robust losses and resistant summaries limit the influence of outliers while changing the estimator's efficiency and uncertainty calculation.}
\noterow{How should it be used?}{For Robust Estimation, begin from explicit inputs, outputs, preconditions, and representation choices.  Specify observation and error models; build the design or covariance structure; fit with a stable solver; inspect residuals and identifiability; quantify uncertainty; test predictions out of sample.}
\noterow{Quantitative lens}{Identify the governing equation, recurrence, probability law, or geometric predicate for Robust Estimation.  Define every variable and scale; then derive a residual, error estimate, or invariant that can be checked independently.}
\noterow{Cost and reliability}{Analyze arithmetic work, storage, data movement, and convergence separately.  Relevant failure pressures include rank deficiency, overfitting, correlated errors, unjustified gaussian assumptions, parameter nonidentifiability, poor mcmc mixing, kernel hyperparameter confounding, and data leakage.  Use scaled tolerances and report both success criteria and diagnostic evidence.}
\noterow{Implementation checklist}{\begin{itemize}
\item Validate dimensions, domains, ordering, and structural assumptions.
\item Guard small divisors, invalid states, overflow, underflow, and nonfinite values.
\item Make mutation, workspace, indexing, and termination behavior explicit.
\item Test a simple exact case, a difficult valid case, a boundary case, and an expected failure.
\end{itemize}}
\end{cornell}
\begin{summarybox}{Section summary}
Robust losses and resistant summaries limit the influence of outliers while changing the estimator's efficiency and uncertainty calculation.  Select this technique only when its assumptions match the data, and accept a result only after an independent residual, error estimate, invariant, or refinement check.
\end{summarybox}

\section{15.8 Markov Chain Monte Carlo}
\begin{cornell}
\noterow{What is the central idea?}{MCMC constructs a correlated chain with the target distribution as its stationary law and requires convergence and effective-sample diagnostics.}
\noterow{How should it be used?}{For Markov Chain Monte Carlo, begin from explicit inputs, outputs, preconditions, and representation choices.  Specify observation and error models; build the design or covariance structure; fit with a stable solver; inspect residuals and identifiability; quantify uncertainty; test predictions out of sample.}
\noterow{Quantitative anchor}{\[
\alpha(x,y)=\min\!\left(1,\frac{\pi(y)q(x\mid y)}{\pi(x)q(y\mid x)}\right)
\]
State the dimensions, scaling, and validity assumptions before using this relation.  Check the resulting residual or invariant rather than trusting an iteration count alone.}
\noterow{Cost and reliability}{Analyze arithmetic work, storage, data movement, and convergence separately.  Relevant failure pressures include rank deficiency, overfitting, correlated errors, unjustified gaussian assumptions, parameter nonidentifiability, poor mcmc mixing, kernel hyperparameter confounding, and data leakage.  Use scaled tolerances and report both success criteria and diagnostic evidence.}
\noterow{Implementation checklist}{\begin{itemize}
\item Validate dimensions, domains, ordering, and structural assumptions.
\item Guard small divisors, invalid states, overflow, underflow, and nonfinite values.
\item Make mutation, workspace, indexing, and termination behavior explicit.
\item Test a simple exact case, a difficult valid case, a boundary case, and an expected failure.
\end{itemize}}
\end{cornell}
\begin{summarybox}{Section summary}
MCMC constructs a correlated chain with the target distribution as its stationary law and requires convergence and effective-sample diagnostics.  Select this technique only when its assumptions match the data, and accept a result only after an independent residual, error estimate, invariant, or refinement check.
\end{summarybox}

\section{15.9 Gaussian Process Regression}
\begin{cornell}
\noterow{What is the central idea?}{A Gaussian process places a covariance-controlled distribution over functions, yielding predictive means and uncertainties after conditioning on data.}
\noterow{How should it be used?}{For Gaussian Process Regression, begin from explicit inputs, outputs, preconditions, and representation choices.  Specify observation and error models; build the design or covariance structure; fit with a stable solver; inspect residuals and identifiability; quantify uncertainty; test predictions out of sample.}
\noterow{Quantitative anchor}{\[
f\sim\mathcal{GP}(m,k)
\]
State the dimensions, scaling, and validity assumptions before using this relation.  Check the resulting residual or invariant rather than trusting an iteration count alone.}
\noterow{Cost and reliability}{Analyze arithmetic work, storage, data movement, and convergence separately.  Relevant failure pressures include rank deficiency, overfitting, correlated errors, unjustified gaussian assumptions, parameter nonidentifiability, poor mcmc mixing, kernel hyperparameter confounding, and data leakage.  Use scaled tolerances and report both success criteria and diagnostic evidence.}
\noterow{Implementation checklist}{\begin{itemize}
\item Validate dimensions, domains, ordering, and structural assumptions.
\item Guard small divisors, invalid states, overflow, underflow, and nonfinite values.
\item Make mutation, workspace, indexing, and termination behavior explicit.
\item Test a simple exact case, a difficult valid case, a boundary case, and an expected failure.
\end{itemize}}
\end{cornell}
\begin{summarybox}{Section summary}
A Gaussian process places a covariance-controlled distribution over functions, yielding predictive means and uncertainties after conditioning on data.  Select this technique only when its assumptions match the data, and accept a result only after an independent residual, error estimate, invariant, or refinement check.
\end{summarybox}

\section*{Method comparison}
\begin{center}
\small
\begin{tabularx}{\linewidth}{>{\bfseries\raggedright\arraybackslash}p{0.22\linewidth}XX}
\toprule
Method or lens & Best use & Main caution \\
\midrule
Linear least squares & Parameters enter linearly & Fast; inspect rank and residual model \\
Nonlinear least squares & Smooth nonlinear parameters & Needs initialization and globalization \\
Gaussian process & Flexible probabilistic function model & Cubic dense cost; kernel choice \\
\bottomrule
\end{tabularx}
\end{center}

\section*{Worked example}
\begin{exambox}{Independent worked example}
Fit $y=a+bx$ to $(0,1),(1,3),(2,5)$.  Centered least squares gives $b=2$ and $a=1$, with zero residuals.  Zero training residual does not prove the line generalizes; uncertainty is not estimable from these exact three points without an error model.
\end{exambox}

\section*{Chapter synthesis}
\begin{summarybox}{Chapter synthesis}
The chapter's methods should be viewed as a decision system rather than a list of routines.  Begin with the mathematical problem and its conditioning, exploit structure, select an algorithm with compatible stability and cost, and verify the computed answer with evidence that is independent of the internal iteration.  The recurring implementation discipline is: specify observation and error models; build the design or covariance structure; fit with a stable solver; inspect residuals and identifiability; quantify uncertainty; test predictions out of sample.
\end{summarybox}

\subsection*{Common mistakes and numerical hazards}
\begin{warnbox}{Common mistakes and numerical hazards}
Rank deficiency, overfitting, correlated errors, unjustified Gaussian assumptions, parameter nonidentifiability, poor MCMC mixing, kernel hyperparameter confounding, and data leakage.  A second recurring mistake is reporting only a computed value: always retain tolerances, iteration counts, residuals, refinement behavior, and relevant structural checks.
\end{warnbox}

\subsection*{Retrieval practice}
\textbf{Questions}
\begin{enumerate}
\item What problem does Least Squares as a Maximum Likelihood Estimator address, and which assumption is easiest to violate?
\item What problem does Fitting Data to a Straight Line address, and which assumption is easiest to violate?
\item What problem does Straight-Line Data with Errors in Both Coordinates address, and which assumption is easiest to violate?
\item What problem does Markov Chain Monte Carlo address, and which assumption is easiest to violate?
\item What problem does Gaussian Process Regression address, and which assumption is easiest to violate?
\item How do conditioning, stability, and the final residual play different roles in this chapter?
\end{enumerate}

\textbf{Brief answers}
\begin{enumerate}
\item Under independent Gaussian errors with known scales, minimizing weighted squared residuals is equivalent to maximizing likelihood. Recheck the section's domain, structure, scaling, and failure diagnostics.
\item Straight-line fitting estimates slope and intercept from centered sums, with uncertainty determined by error assumptions and leverage. Recheck the section's domain, structure, scaling, and failure diagnostics.
\item When both coordinates are uncertain, ordinary vertical residuals are inadequate and the likelihood must account for two-axis error geometry. Recheck the section's domain, structure, scaling, and failure diagnostics.
\item MCMC constructs a correlated chain with the target distribution as its stationary law and requires convergence and effective-sample diagnostics. Recheck the section's domain, structure, scaling, and failure diagnostics.
\item A Gaussian process places a covariance-controlled distribution over functions, yielding predictive means and uncertainties after conditioning on data. Recheck the section's domain, structure, scaling, and failure diagnostics.
\item Conditioning describes sensitivity of the mathematical problem, stability describes error propagation by the algorithm, and the residual measures satisfaction of the computed equations; none is a universal substitute for the others.
\end{enumerate}

\subsection*{Mastery checklist}
\begin{itemize}
\item[$\square$] I can explain and select Least Squares as a Maximum Likelihood Estimator without relying on a memorized code listing.
\item[$\square$] I can explain and select Fitting Data to a Straight Line without relying on a memorized code listing.
\item[$\square$] I can explain and select Straight-Line Data with Errors in Both Coordinates without relying on a memorized code listing.
\item[$\square$] I can explain and select General Linear Least Squares without relying on a memorized code listing.
\item[$\square$] I can explain and select Nonlinear Models without relying on a memorized code listing.
\item[$\square$] I can state a scaled stopping rule and an independent verification check.
\item[$\square$] I can identify at least one conditioning risk and one algorithmic stability risk.
\item[$\square$] I can estimate time, storage, and data-movement costs at the appropriate level.
\end{itemize}

\end{document}

